\section{AMC 10 10B 2024}
        \begin{enumerate}
        	\item (\textit{AMC 10 10B 2024}) In a long line of people arranged left to right, the 1013th person from the left is also the 1010th person from the right. How many people are in the line?


 $\textbf{(A) } 2021 \qquad\textbf{(B) } 2022 \qquad\textbf{(C) } 2023 \qquad\textbf{(D) } 2024 \qquad\textbf{(E) } 2025$


	\item (\textit{AMC 10 10B 2024}) What is $10! - 7! \cdot 6!?$
$\textbf{(A) } -120 \qquad\textbf{(B) } 0 \qquad\textbf{(C) } 120 \qquad\textbf{(D) } 600 \qquad\textbf{(E) } 720$


	\item (\textit{AMC 10 10B 2024}) For how many integer values of $x$ is $|2x| \leq 7 \pi$?


 $\textbf{(A) } 16 \qquad\textbf{(B) } 17 \qquad\textbf{(C) } 19 \qquad\textbf{(D) } 20 \qquad\textbf{(E) } 21$


	\item (\textit{AMC 10 10B 2024}) Balls numbered $1, 2, 3, \dots$ are deposited in $5$ bins, labeled $A$, $B$, $C$, $D$, and $E$ using the following procedure. Ball $1$ is deposited in bin $A$, and balls $2$ and $3$ are deposited in bin $B$. The next $3$ balls are deposited in bin $C$, the next $4$ in bin $D$, and so on, cycling back to bin $A$ after balls are deposited in bin $E$. (For example, balls numbered $22, 23, \dots, 28$ are deposited in bin $B$ at step $7$ of this process.) In which bin is ball $2024$ deposited?


 $\textbf{(A) } A \qquad\textbf{(B) } B \qquad\textbf{(C) } C \qquad\textbf{(D) } D \qquad\textbf{(E) } E$


	\item (\textit{AMC 10 10B 2024}) In the following expression, Melanie changed some of the plus signs to minus signs:

 \[1+3+5+7+\cdots+97+99\]
When the new expression was evaluated, it was negative. What is the least number of plus signs that Melanie could have changed to minus signs?


 $\textbf{(A) } 14 \qquad\textbf{(B) } 15 \qquad\textbf{(C) } 16 \qquad\textbf{(D) } 17 \qquad\textbf{(E) } 18$


	\item (\textit{AMC 10 10B 2024}) A rectangle has integer side lengths and an area of $2024$. What is the least possible perimeter of the rectangle?


 $\textbf{(A) } 160 \qquad\textbf{(B) } 180 \qquad\textbf{(C) } 222 \qquad\textbf{(D) } 228 \qquad\textbf{(E) } 390$


	\item (\textit{AMC 10 10B 2024}) What is the remainder when $7^{2024}+7^{2025}+7^{2026}$ is divided by $19$?


 $\textbf{(A) } 0 \qquad\textbf{(B) } 1 \qquad\textbf{(C) } 7 \qquad\textbf{(D) } 11 \qquad\textbf{(E) } 18$


	\item (\textit{AMC 10 10B 2024}) Let $N$ be the product of all the positive integer divisors of $42$. What is the units digit
of $N$?


 $\textbf{(A) } 0\qquad\textbf{(B) } 2\qquad\textbf{(C) } 4\qquad\textbf{(D) } 6\qquad\textbf{(E) } 8$


	\item (\textit{AMC 10 10B 2024}) Real numbers $a, b,$ and $c$ have arithmetic mean $0$. The arithmetic mean of $a^2, b^2,$ and $c^2$ is $10$. What is the arithmetic mean of $ab, ac,$ and $bc$?


 $\textbf{(A) } -5 \qquad\textbf{(B) } -\dfrac{10}{3} \qquad\textbf{(C) } -\dfrac{10}{9} \qquad\textbf{(D) } 0 \qquad\textbf{(E) } \dfrac{10}{9}$


	\item (\textit{AMC 10 10B 2024}) Quadrilateral $ABCD$ is a parallelogram, and $E$ is the midpoint of the side $\overline{AD}$. Let $F$ be the intersection of lines $EB$ and $AC$. What is the ratio of the area of quadrilateral $CDEF$ to the area of $\triangle CFB$?


 $\textbf{(A) } 5:4 \qquad\textbf{(B) } 4:3 \qquad\textbf{(C) } 3:2 \qquad\textbf{(D) } 5:3 \qquad\textbf{(E) } 2:1$


	\item (\textit{AMC 10 10B 2024}) In the figure below $WXYZ$ is a rectangle with $WX=4$ and $WZ=8$. Point $M$ lies on $\overline{XY}$, point $A$ lies on $\overline{YZ}$, and $\angle WMA$ is a right angle. The areas of $\triangle WXM$ and $\triangle W A Z$ are equal. What is the area of $\triangle WMA$?


 
                    \begin{figure}[H]
                        \centering
                        \includegraphics[width=0.5\textwidth]{Problem/AMC_10_10B_2024_11_Problem_Figure_1.png}
                    \end{figure}
                    
$\textbf{(A) }13 \qquad \textbf{(B) }14 \qquad \textbf{(C) }15 \qquad \textbf{(D) }16 \qquad \textbf{(E) }17 \qquad$


	\item (\textit{AMC 10 10B 2024}) A group of $100$ students from different countries meet at a mathematics competition.
Each student speaks the same number of languages, and, for every pair of
students $A$ and $B$, student $A$ speaks some language that student $B$ does not speak,
and student $B$ speaks some language that student $A$ does not speak. What is the
least possible total number of languages spoken by all the students?


 $\textbf{(A) } 9 \qquad\textbf{(B) } 10 \qquad\textbf{(C) } 12 \qquad\textbf{(D) } 51 \qquad\textbf{(E) } 100$


	\item (\textit{AMC 10 10B 2024}) Positive integers $x$ and $y$ satisfy the equation $\sqrt{x} + \sqrt{y} = \sqrt{1183}$. What is the minimum possible value of $x+y$?


 $\textbf{(A) } 585 \qquad\textbf{(B) } 595 \qquad\textbf{(C) } 623 \qquad\textbf{(D) } 700 \qquad\textbf{(E) } 791$


	\item (\textit{AMC 10 10B 2024}) A dartboard is the region $B$ in the coordinate plane consisting of points $(x,y)$ such
that $|x| + |y| \le 8$ . A target $T$ is the region where $(x^2 + y^2 - 25)^2 \le 49.$ A dart is
thrown and lands at a random point in $B$. The probability that the dart lands in $T$ can be
expressed as $\frac{m}{n} \cdot \pi,$ where $m$ and $n$ are relatively prime positive integers. What
is $m + n?$
$\textbf{(A) } 39 \qquad\textbf{(B) } 71 \qquad\textbf{(C) } 73 \qquad\textbf{(D) } 75 \qquad\textbf{(E) } 135$


	\item (\textit{AMC 10 10B 2024}) A list of $9$ real numbers consists of $1$, $2.2$, $3.2$, $5.2$, $6.2$, and $7$, as well as $x$, $y$ , and $z$ with $x$ $\le$ $y$ $\le$ $z$. The range of the list is $7$, and the mean and the median are both positive integers. How many ordered triples ($x$, $y$, $z$) are possible?


 $\textbf{(A) } 1 \qquad \textbf{(B) } 2 \qquad \textbf{(C) } 3 \qquad \textbf{(D) } 4 \qquad \textbf{(E) } \text{infinitely many}$


	\item (\textit{AMC 10 10B 2024}) Jerry likes to play with numbers. One day, he wrote all the integers from $1$ to $2024$ on the whiteboard. Then he repeatedly chose four numbers on the whiteboard, erased them, and replaced them by either their sum or their product. (For example, Jerry's first step might have been to erase $1$, $2$, $3$, and $5$, and then write either $11$, their sum, or $30$, their product, on the whiteboard.) After repeatedly performing this operation, Jerry noticed that all the remaining numbers on the whiteboard were odd. What is the maximum possible number of integers on the whiteboard at that time?


 $\textbf{(A) } 1010 \qquad \textbf{(B) } 1011 \qquad \textbf{(C) } 1012 \qquad \textbf{(D) } 1013 \qquad \textbf{(E) } 1014$


	\item (\textit{AMC 10 10B 2024}) In a race among 5 snails, there is at most one tie, but that tie can involve any number
of snails. For example, the result of the race might be that Dazzler is first; Abby,
Cyrus, and Elroy are tied for second; and Bruna is fifth. How many different results
of the race are possible?


 $\textbf{(A) } 180 \qquad\textbf{(B) } 361 \qquad\textbf{(C) } 420 \qquad\textbf{(D) } 431 \qquad\textbf{(E) } 720$


	\item (\textit{AMC 10 10B 2024}) How many different remainders can result when the $100$th power of an integer is
divided by $125$?


 $\textbf{(A) } 1 \qquad\textbf{(B) } 2 \qquad\textbf{(C) } 5 \qquad\textbf{(D) } 25 \qquad\textbf{(E) } 125$


	\item (\textit{AMC 10 10B 2024}) In the following table, each question mark is to be replaced by "Possible" or "Not
Possible" to indicate whether a nonvertical line with the given slope can contain the
given number of lattice points (points both of whose coordinates are integers). How
many of the 12 entries will be "Possible"?


 
                    \begin{figure}[H]
                        \centering
                        \includegraphics[width=0.5\textwidth]{Problem/AMC_10_10B_2024_19_Problem_Figure_1.png}
                    \end{figure}
                    
$\textbf{(A) } 4 \qquad     \textbf{(B) } 5 \qquad \textbf{(C) } 6    \qquad \textbf{(D) } 7      \qquad \textbf{(E) } 9$


	\item (\textit{AMC 10 10B 2024}) Three different pairs of shoes are placed in a row so that no left shoe is next to a
right shoe from a different pair. In how many ways can these six shoes be lined up?


 $\textbf{(A) } 60 \qquad\textbf{(B) } 72 \qquad\textbf{(C) } 90 \qquad\textbf{(D) } 108 \qquad\textbf{(E) } 120$


	\item (\textit{AMC 10 10B 2024}) Two straight pipes (circular cylinders), with radii $1$ and $\frac{1}{4}$, lie parallel and in contact on a flat floor. The figure below shows a head-on view. What is the sum of the possible radii of a third parallel pipe lying on the same floor and in contact with both?


 
                    \begin{figure}[H]
                        \centering
                        \includegraphics[width=0.5\textwidth]{Problem/AMC_10_10B_2024_21_Problem_Figure_1.png}
                    \end{figure}
                    
$\textbf{(A)}~\frac{1}{9} \qquad\textbf{(B)}~1 \qquad\textbf{(C)}~\frac{10}{9} \qquad\textbf{(D)}~\frac{11}{9} \qquad\textbf{(E)}~\frac{19}{9}$


	\item (\textit{AMC 10 10B 2024}) A group of $16$ people will be partitioned into $4$ indistinguishable $4$-person
committees. Each committee will have one chairperson and one secretary. The
number of different ways to make these assignments can be written
as ${3^r}M$, where $r$ and $M$ are positive integers and $M$ is not divisible by $3$. What is $r$?


 $\textbf{(A) } 5 \qquad\textbf{(B) } 6 \qquad\textbf{(C) } 7 \qquad\textbf{(D) } 8 \qquad\textbf{(E) } 9$


	\item (\textit{AMC 10 10B 2024}) The Fibonacci numbers are defined by $F_1 = 1, F_2 = 1,$ and $F_n = F_{n-1} + F_{n-2}$ for $n \geq 3.$ What is 
 \[{\frac{F_2}{F_1}} + {\frac{F_4}{F_2}} + {\frac{F_6}{F_3}} + ... + {\frac{F_{20}}{F_{10}}}?\]
$\textbf{(A) } 318 \qquad\textbf{(B) } 319 \qquad\textbf{(C) } 320 \qquad\textbf{(D) } 321 \qquad\textbf{(E) } 322$


	\item (\textit{AMC 10 10B 2024}) Let

 \[P(m)=\frac{m}{2}+\frac{m^2}{4}+\frac{m^4}{8}+\frac{m^8}{8}\]
How many of the values $P(2022)$, $P(2023)$, $P(2024)$, and $P(2025)$ are integers?


 $\textbf{(A) } 0 \qquad\textbf{(B) } 1 \qquad\textbf{(C) } 2 \qquad\textbf{(D) } 3 \qquad\textbf{(E) } 4$


	\item (\textit{AMC 10 10B 2024}) Each of $27$ bricks (right rectangular prisms) has dimensions $a \times b \times c$, where $a$, $b$, and $c$ are pairwise relatively prime positive integers. These bricks are arranged to form a $3 \times 3 \times 3$ block, as shown on the left below. A $28$th brick with the same dimensions is introduced, and these bricks are reconfigured into a $2 \times 2 \times 7$ block, shown on the right. The new block is $1$ unit taller, $1$ unit wider, and $1$ unit deeper than the old one. What is $a + b + c$?


 
                    \begin{figure}[H]
                        \centering
                        \includegraphics[width=0.5\textwidth]{Problem/AMC_10_10B_2024_25_Problem_Figure_1.png}
                    \end{figure}
                    
$\textbf{(A) } 88 \qquad\textbf{(B) } 89 \qquad\textbf{(C) } 90 \qquad\textbf{(D) } 91 \qquad\textbf{(E) } 92$


\end{enumerate}